\documentclass{beamer}
\usetheme{Madrid}
\usecolortheme{default}
\usepackage[utf8]{inputenc}
\usepackage{amsmath}
\usepackage{amsfonts}
\usepackage{amssymb}
\usepackage{tikz}

\title{CRISC Domain C - 100 Very Hard Questions}
\subtitle{Risk Response and Reporting - Expert Level}
\author{CRISC Mastery Series}
\date{\today}

\begin{document}
	
	\begin{frame}
		\titlepage
	\end{frame}
	
	\begin{frame}{Outline}
		\tableofcontents
	\end{frame}
	
	% ============================================================
	% SECTION 1: RISK RESPONSE STRATEGIES - HARD QUESTIONS
	% ============================================================
	
	\section{1. Risk Response Strategies}
	
	\begin{frame}{Q1: Mitigate vs Transfer Trap}
		\begin{block}{Question}
			An organization outsources its payroll processing to a third-party vendor. The contract includes indemnification clauses requiring the vendor to repay any financial losses caused by vendor negligence. The organization's risk manager classifies this as risk transfer. Is this correct?
		\end{block}
		\begin{enumerate}[(A)]
			\item Yes, because the contract shifts liability to the vendor
			\item No, because the organization retains accountability for payroll data
			\item Yes, because indemnification clauses constitute transfer
			\item No, because outsourcing always constitutes avoidance
		\end{enumerate}
	\end{frame}
	
	\begin{frame}{Q1: Explanation}
		\begin{block}{Correct Answer: B}
			Outsourcing with indemnification is a combination of mitigation (reducing likelihood through vendor expertise) and partial transfer (financial indemnity). However, the organization \textbf{retains ultimate accountability} for payroll data and compliance. True transfer requires shifting both financial impact AND responsibility to another party (e.g., insurance). The reputational and regulatory accountability remains with the organization.
		\end{block}
		\begin{alertblock}{Critical Trap}
			Outsourcing = Mitigation + Partial Transfer, NOT pure Transfer. Accountability never transfers.
		\end{alertblock}
	\end{frame}
	
	\begin{frame}{Q2: Cost-Benefit Decision Logic}
		\begin{block}{Question}
			A risk has an inherent risk exposure of \$500,000 annually. The cost to implement controls to reduce this exposure to \$50,000 is \$200,000 annually. However, the organization's risk appetite is \$75,000. What is the MOST appropriate action?
		\end{block}
		\begin{enumerate}[(A)]
			\item Accept the risk because residual risk is below \$500,000
			\item Implement the controls because residual risk (\$50,000) is below risk appetite (\$75,000)
			\item Transfer the risk because it exceeds risk appetite
			\item Avoid the risk entirely because it exceeds risk appetite
		\end{enumerate}
	\end{frame}
	
	\begin{frame}{Q2: Explanation}
		\begin{block}{Correct Answer: B}
			The cost-benefit analysis: \$200,000 control cost < \$450,000 risk reduction (\$500,000 - \$50,000). The residual risk (\$50,000) falls below the risk appetite (\$75,000). This is the textbook scenario where mitigation is justified both economically and strategically.
		\end{block}
		\begin{block}{Why Others Are Wrong}
			\begin{itemize}
				\item \textbf{A:} Acceptance leaves risk above appetite (\$500,000 > \$75,000)
				\item \textbf{C:} Transfer may not be cost-effective if controls work
				\item \textbf{D:} Avoidance is extreme when controls can achieve acceptable risk
			\end{itemize}
		\end{block}
	\end{frame}
	
	\begin{frame}{Q3: Risk Avoidance Irreversibility}
		\begin{block}{Question}
			A company decides to shut down its e-commerce platform due to persistent DDoS attacks. The CEO argues this is risk mitigation because they are reducing the attack surface. The risk manager disagrees. Who is correct and why?
		\end{block}
		\begin{enumerate}[(A)]
			\item CEO is correct because reducing attack surface is mitigation
			\item Risk manager is correct because termination of activity is avoidance
			\item Both are correct - it can be classified either way
			\item Neither is correct - this is risk acceptance
		\end{enumerate}
	\end{frame}
	
	\begin{frame}{Q3: Explanation}
		\begin{block}{Correct Answer: B}
			\textbf{Avoidance} = terminating the activity that gives rise to risk entirely. Shutting down the e-commerce platform eliminates the DDoS risk because the platform no longer exists. This is NOT mitigation (which would involve implementing DDoS protection while keeping the platform operational). Avoidance is irreversible and requires senior management approval.
		\end{block}
		\begin{alertblock}{Critical Trap}
			Shutting down a business function = Avoidance, NOT Mitigation.
			Mitigation = Keep the function, reduce the risk.
		\end{alertblock}
	\end{frame}
	
	\begin{frame}{Q4: Transfer with Residual Reputational Risk}
		\begin{block}{Question}
			An organization purchases cyber insurance covering \$10 million in breach costs. After a breach occurs and the insurance pays out, what risk remains with the organization?
		\end{block}
		\begin{enumerate}[(A)]
			\item No risk remains because insurance covers all costs
			\item Only the deductible amount remains as risk
			\item Reputational damage, customer trust loss, and regulatory accountability
			\item Only legal fees remain as risk
		\end{enumerate}
	\end{frame}
	
	\begin{frame}{Q4: Explanation}
		\begin{block}{Correct Answer: C}
			Insurance transfers \textbf{financial} risk but NOT reputational, strategic, or accountability risk. After a breach:
			\begin{itemize}
				\item Customers may leave (revenue loss)
				\item Brand reputation is damaged
				\item Regulators still hold the organization accountable
				\item Lawsuits may still name the organization
			\end{itemize}
		\end{block}
		\begin{alertblock}{Critical Trap}
			Transfer = Financial only. Reputational and accountability risk ALWAYS remains.
		\end{alertblock}
	\end{frame}
	
	\begin{frame}{Q5: Acceptance with Unknown Risk}
		\begin{block}{Question}
			A business unit head accepts a risk without understanding its potential impact because "it hasn't happened yet." This is an example of:
		\end{block}
		\begin{enumerate}[(A)]
			\item Informed risk acceptance
			\item Uninformed risk acceptance (negligence)
			\item Risk avoidance
			\item Risk transfer
		\end{enumerate}
	\end{frame}
	
	\begin{frame}{Q5: Explanation}
		\begin{block}{Correct Answer: B}
			Risk acceptance must be \textbf{informed}. Accepting risk without understanding impact is negligence and could lead to legal liability. Proper acceptance requires:
			\begin{itemize}
				\item Understanding the risk (likelihood and impact)
				\item Understanding the consequences of acceptance
				\item Formal documentation and sign-off
				\item Periodic review
			\end{itemize}
		\end{block}
		\begin{alertblock}{Critical Trap}
			Acceptance without understanding = Negligence, NOT valid acceptance.
		\end{alertblock}
	\end{frame}
	
	\begin{frame}{Q6: Mitigation Cost Exceeds Asset Value}
		\begin{block}{Question}
			A \$5,000 laptop contains sensitive data. The cost to implement full-disk encryption and endpoint protection is \$8,000. What is the BEST risk response?
		\end{block}
		\begin{enumerate}[(A)]
			\item Implement the controls anyway because data is sensitive
			\item Accept the risk because control cost exceeds asset value
			\item Transfer the risk through insurance
			\item Avoid the risk by not using laptops
		\end{enumerate}
	\end{frame}
	
	\begin{frame}{Q6: Explanation}
		\begin{block}{Correct Answer: B}
			The fundamental principle: \textbf{cost of control should not exceed the value of the asset}. Since \$8,000 > \$5,000, the control is not economically justified. Accept the risk (with documentation) or consider lower-cost controls. However, note that if the data value exceeds the laptop value, the equation changes.
		\end{block}
		\begin{alertblock}{Critical Trap}
			Cost of control vs. Asset value is the primary decision factor.
			\$8,000 control on \$5,000 asset = Accept the risk.
		\end{alertblock}
	\end{frame}
	
	\begin{frame}{Q7: Risk Response Selection with Multiple Factors}
		\begin{block}{Question}
			A company faces a risk of \$2M from a data breach. Options:
			\begin{itemize}
				\item Option A: Implement controls costing \$500K, reducing risk to \$200K
				\item Option B: Purchase insurance costing \$300K, covering \$1.5M
				\item Option C: Outsource to vendor costing \$400K, transferring operational risk
			\end{itemize}
			The company's risk appetite is \$250K. What is the BEST response?
		\end{block}
		\begin{enumerate}[(A)]
			\item Option A only
			\item Option B only
			\item Option C only
			\item Option A and B combined
		\end{enumerate}
	\end{frame}
	
	\begin{frame}{Q7: Explanation}
		\begin{block}{Correct Answer: D}
			\textbf{Option A + B Combined:}
			\begin{itemize}
				\item Controls reduce risk to \$200K (below appetite of \$250K)
				\item Insurance covers remaining \$200K financial exposure
				\item Total cost: \$800K for \$2M risk reduction = strong ROI
				\item Residual risk after both: \$0 financial, but reputational risk remains
			\end{itemize}
		\end{block}
		\begin{block}{Why Not Others}
			\begin{itemize}
				\item A alone: Residual \$200K is below appetite but still significant
				\item B alone: Doesn't reduce likelihood, only financial impact
				\item C alone: Doesn't address the core risk, only shifts operations
			\end{itemize}
		\end{block}
	\end{frame}
	
	\begin{frame}{Q8: Avoidance vs. Acceptance Trap}
		\begin{block}{Question}
			A company decides not to pursue a new market because the regulatory risk is too high. This is:
		\end{block}
		\begin{enumerate}[(A)]
			\item Risk acceptance because they are accepting the opportunity cost
			\item Risk avoidance because they are not entering the market
			\item Risk mitigation because they are reducing exposure
			\item Risk transfer because they are shifting risk to competitors
		\end{enumerate}
	\end{frame}
	
	\begin{frame}{Q8: Explanation}
		\begin{block}{Correct Answer: B}
			\textbf{Avoidance} = Not engaging in the activity that creates the risk. By not entering the market, the company avoids the regulatory risk entirely. This is a valid avoidance strategy.
		\end{block}
		\begin{alertblock}{Critical Trap}
			Not doing something = Avoidance.
			Doing something and accepting the consequences = Acceptance.
		\end{alertblock}
	\end{frame}
	
	\begin{frame}{Q9: Transfer Limitations}
		\begin{block}{Question}
			Which risk CANNOT be effectively transferred?
		\end{block}
		\begin{enumerate}[(A)]
			\item Financial loss from a data breach
			\item Regulatory accountability for data protection
			\item Property damage from a natural disaster
			\item Liability from a product defect
		\end{enumerate}
	\end{frame}
	
	\begin{frame}{Q9: Explanation}
		\begin{block}{Correct Answer: B}
			\textbf{Regulatory accountability cannot be transferred.} While you can transfer financial liability through insurance or contracts, regulators hold the \textbf{original organization} accountable for compliance. You cannot "outsource" your regulatory responsibility.
		\end{block}
		\begin{block}{What Can Be Transferred}
			\begin{itemize}
				\item Financial losses (insurance)
				\item Operational risk (outsourcing)
				\item Liability (contracts with indemnification)
			\end{itemize}
		\end{block}
		\begin{alertblock}{Critical Trap}
			Regulatory accountability = Non-transferable.
			You can never transfer your legal responsibility to a regulator.
		\end{alertblock}
	\end{frame}
	
	\begin{frame}{Q10: Mitigation vs. Compensating Controls}
		\begin{block}{Question}
			A legacy system cannot be patched due to vendor restrictions. The security team implements network segmentation and enhanced monitoring. This is:
		\end{block}
		\begin{enumerate}[(A)]
			\item Risk avoidance because the vulnerability is eliminated
			\item Risk mitigation through compensating controls
			\item Risk acceptance because the system remains vulnerable
			\item Risk transfer because monitoring is outsourced
		\end{enumerate}
	\end{frame}
	
	\begin{frame}{Q10: Explanation}
		\begin{block}{Correct Answer: B}
			\textbf{Compensating controls are a form of mitigation.} When primary controls cannot be implemented, compensating controls reduce the likelihood or impact of exploitation. Network segmentation reduces lateral movement; enhanced monitoring increases detection.
		\end{block}
		\begin{block}{Why Not Others}
			\begin{itemize}
				\item A: The vulnerability is NOT eliminated (system is still unpatched)
				\item C: The organization is NOT passively accepting (they implemented controls)
				\item D: Monitoring being outsourced doesn't make it transfer
			\end{itemize}
		\end{block}
	\end{frame}
	
	% ============================================================
	% SECTION 2: RISK AND CONTROL OWNERSHIP - HARD QUESTIONS
	% ============================================================
	
	\section{2. Risk and Control Ownership}
	
	\begin{frame}{Q11: Risk Owner vs. Control Owner Conflict}
		\begin{block}{Question}
			The risk owner for a critical application wants to accept a high-severity vulnerability because remediation costs \$500K. The control owner insists on remediation due to regulatory requirements. Who has the authority to make the final decision?
		\end{block}
		\begin{enumerate}[(A)]
			\item Risk owner, because they own the risk
			\item Control owner, because they own the control
			\item Senior management, because it involves regulatory risk
			\item Both must agree for any decision to be made
		\end{enumerate}
	\end{frame}
	
	\begin{frame}{Q11: Explanation}
		\begin{block}{Correct Answer: C}
			\textbf{Senior management} must make the decision because:
			\begin{itemize}
				\item Regulatory risk affects the entire organization, not just one business unit
				\item Risk acceptance above risk appetite requires executive approval
				\item Legal/regulatory obligations override cost considerations
				\item Risk owner cannot unilaterally accept regulatory non-compliance
			\end{itemize}
		\end{block}
		\begin{alertblock}{Critical Trap}
			Risk owner not equal to Final authority on ALL risk decisions.
			Regulatory/legal risk = Senior management decision.
		\end{alertblock}
	\end{frame}
	
	\begin{frame}{Q12: Multiple Risk Owners}
		\begin{block}{Question}
			A risk affects three business units: Marketing, Sales, and Operations. Each BU head wants to be the risk owner. What is the BEST approach?
		\end{block}
		\begin{enumerate}[(A)]
			\item Appoint all three as co-owners
			\item Appoint the BU with the largest budget as owner
			\item Appoint a single owner with authority across all BUs
			\item Escalate to the board for decision
		\end{enumerate}
	\end{frame}
	
	\begin{frame}{Q12: Explanation}
		\begin{block}{Correct Answer: C}
			Risk ownership requires \textbf{single-point accountability}. A risk owner must have:
			\begin{itemize}
				\item Authority to make decisions across all affected areas
				\item Budget control to implement risk responses
				\item Accountability for outcomes
			\end{itemize}
			Co-ownership dilutes accountability. The owner should be the executive with cross-functional authority (e.g., CRO, COO).
		\end{block}
		\begin{alertblock}{Critical Trap}
			Multiple owners = No owner.
			Single point of accountability is required.
		\end{alertblock}
	\end{frame}
	
	\begin{frame}{Q13: Control Owner Without Authority}
		\begin{block}{Question}
			A control owner identifies that a critical control is ineffective but has no budget to fix it. What should they do FIRST?
		\end{block}
		\begin{enumerate}[(A)]
			\item Document the finding and escalate to the risk owner
			\item Implement a compensating control anyway
			\item Accept the risk and move on
			\item Wait for the next audit cycle to report
		\end{enumerate}
	\end{frame}
	
	\begin{frame}{Q13: Explanation}
		\begin{block}{Correct Answer: A}
			The control owner must \textbf{escalate to the risk owner} who has the budget authority. The control owner is responsible for the control's effectiveness but may not have budget authority. The risk owner is accountable for ensuring the risk is treated.
		\end{block}
		\begin{block}{Why Not Others}
			\begin{itemize}
				\item B: Compensating controls require risk owner approval
				\item C: Control owner cannot unilaterally accept risk
				\item D: Delaying reporting increases exposure
			\end{itemize}
		\end{block}
	\end{frame}
	
	\begin{frame}{Q14: Risk Owner Outside IT}
		\begin{block}{Question}
			The risk owner for a cybersecurity risk is the VP of Marketing. IT disagrees with this appointment. Who is correct?
		\end{block}
		\begin{enumerate}[(A)]
			\item IT is correct - cybersecurity risks must be owned by IT
			\item VP of Marketing is correct if the risk impacts marketing objectives
			\item Neither is correct - the CISO should own all cyber risks
			\item The board should decide who owns the risk
		\end{enumerate}
	\end{frame}
	
	\begin{frame}{Q14: Explanation}
		\begin{block}{Correct Answer: B}
			\textbf{Business risk is owned by business leaders.} If a cyber risk (e.g., data breach of customer marketing data) impacts marketing objectives, the VP of Marketing is the appropriate owner. IT is an enabler, not the owner of business risk.
		\end{block}
		\begin{alertblock}{Critical Trap}
			IT does NOT own business risk.
			The business leader who owns the objective owns the risk.
		\end{alertblock}
	\end{frame}
	
	\begin{frame}{Q15: Delegation of Risk Ownership}
		\begin{block}{Question}
			A risk owner goes on extended leave. They delegate risk ownership to a subordinate. What must happen?
		\end{block}
		\begin{enumerate}[(A)]
			\item Nothing - delegation is automatic
			\item The delegation must be documented and communicated to stakeholders
			\item The risk should be transferred to another department
			\item The risk should be accepted until the owner returns
		\end{enumerate}
	\end{frame}
	
	\begin{frame}{Q15: Explanation}
		\begin{block}{Correct Answer: B}
			Delegation must be \textbf{formal, documented, and communicated.} The delegate must:
			\begin{itemize}
				\item Have the authority and budget to make decisions
				\item Understand the risk and current status
				\item Be known to all stakeholders
				\item Have clear boundaries on what they can/cannot approve
			\end{itemize}
		\end{block}
		\begin{alertblock}{Critical Trap}
			Undocumented delegation = No accountability.
			Always document delegations formally.
		\end{alertblock}
	\end{frame}
	
	% ============================================================
	% SECTION 3: THIRD-PARTY RISK MANAGEMENT - HARD QUESTIONS
	% ============================================================
	
	\section{3. Third-Party Risk Management}
	
	\begin{frame}{Q16: Vendor Breach Accountability}
		\begin{block}{Question}
			A SaaS vendor suffers a breach exposing your customer data. The vendor's contract includes a \$5M indemnification clause. Who is held accountable by regulators?
		\end{block}
		\begin{enumerate}[(A)]
			\item The SaaS vendor
			\item Your organization
			\item Both organizations equally
			\item The indemnification clause shifts all accountability
		\end{enumerate}
	\end{frame}
	
	\begin{frame}{Q16: Explanation}
		\begin{block}{Correct Answer: B}
			\textbf{Your organization is held accountable by regulators.} You cannot outsource regulatory accountability. The indemnification clause provides \textbf{financial compensation} but does not shift legal responsibility.
		\end{block}
		\begin{alertblock}{Critical Trap}
			Regulators hold YOU accountable for vendor breaches.
			Indemnification = Financial only, NOT accountability.
		\end{alertblock}
	\end{frame}
	
	\begin{frame}{Q17: Fourth-Party Risk}
		\begin{block}{Question}
			Your critical SaaS vendor uses a subcontractor for data storage. The subcontractor suffers a breach. Which statement is TRUE?
		\end{block}
		\begin{enumerate}[(A)]
			\item Only the subcontractor is liable
			\item Your organization has fourth-party risk exposure
			\item The SaaS vendor bears all responsibility
			\item Regulatory liability ends at your organization
		\end{enumerate}
	\end{frame}
	
	\begin{frame}{Q17: Explanation}
		\begin{block}{Correct Answer: B}
			\textbf{Fourth-party risk} = risk from your vendor's vendors. You have exposure because:
			\begin{itemize}
				\item Your data is held by the subcontractor
				\item Your SaaS vendor may not have adequate oversight
				\item Regulators still hold you accountable
				\item You may not have contractual privity with the subcontractor
			\end{itemize}
		\end{block}
		\begin{alertblock}{Critical Trap}
			Fourth-party risk = Your vendor's vendor.
			You are still exposed even without a direct contract.
		\end{alertblock}
	\end{frame}
	
	\begin{frame}{Q18: Vendor Assessment Timing}
		\begin{block}{Question}
			When is the BEST time to conduct a vendor security assessment?
		\end{block}
		\begin{enumerate}[(A)]
			\item Before signing the contract
			\item After the contract is signed but before go-live
			\item After go-live in the first 30 days
			\item Annually after go-live
		\end{enumerate}
	\end{frame}
	
	\begin{frame}{Q18: Explanation}
		\begin{block}{Correct Answer: A}
			The assessment should occur \textbf{before signing the contract} so that:
			\begin{itemize}
				\item You can negotiate security requirements into the contract
				\item You can make an informed decision about the vendor
				\item You can identify deal-breakers before commitment
				\item You have leverage to require improvements
			\end{itemize}
		\end{block}
		\begin{alertblock}{Critical Trap}
			Assessment AFTER contract = No leverage.
			Assessment BEFORE contract = Negotiating power.
		\end{alertblock}
	\end{frame}
	
	\begin{frame}{Q19: SOC 2 Report Limitations}
		\begin{block}{Question}
			A vendor provides a SOC 2 Type I report. What does this NOT provide?
		\end{block}
		\begin{enumerate}[(A)]
			\item Evidence of control design at a point in time
			\item Evidence of control operating effectiveness over a period
			\item Evidence of management's description of controls
			\item Evidence of the auditor's opinion
		\end{enumerate}
	\end{frame}
	
	\begin{frame}{Q19: Explanation}
		\begin{block}{Correct Answer: B}
			\textbf{SOC 2 Type I} covers control \textbf{design} at a specific point in time. \textbf{SOC 2 Type II} covers control \textbf{operating effectiveness} over a period (typically 6-12 months). Type I does NOT prove the controls operated effectively.
		\end{block}
		\begin{alertblock}{Critical Trap}
			Type I = Design only (snapshot).
			Type II = Design + Operating effectiveness (period).
		\end{alertblock}
	\end{frame}
	
	\begin{frame}{Q20: Vendor Contract Must-Haves}
		\begin{block}{Question}
			Which clause is MOST critical in a vendor contract for data protection?
		\end{block}
		\begin{enumerate}[(A)]
			\item Service level agreement with financial penalties
			\item Right to audit and security requirements
			\item Indemnification for all losses
			\item Non-disclosure agreement
		\end{enumerate}
	\end{frame}
	
	\begin{frame}{Q20: Explanation}
		\begin{block}{Correct Answer: B}
			\textbf{Right to audit + security requirements} is most critical because:
			\begin{itemize}
				\item It provides ongoing verification of compliance
				\item It establishes minimum security standards
				\item It gives you recourse if standards are not met
				\item Financial penalties (A) don't prevent breaches
				\item Indemnification (C) doesn't prevent breaches
			\end{itemize}
		\end{block}
		\begin{alertblock}{Critical Trap}
			Financial penalties don't prevent breaches.
			Right to audit + security requirements = Prevention.
		\end{alertblock}
	\end{frame}
	
	\begin{frame}{Q21: Exception Management Duration}
		\begin{block}{Question}
			An exception to policy is approved for a critical business function. How long can it remain open?
		\end{block}
		\begin{enumerate}[(A)]
			\item Until the next audit cycle
			\item Maximum 1 year, reviewed at least annually
			\item Indefinitely if the business need continues
			\item Maximum 30 days
		\end{enumerate}
	\end{frame}
	
	\begin{frame}{Q21: Explanation}
		\begin{block}{Correct Answer: B}
			Exceptions should be reviewed \textbf{at least annually}. The maximum duration should be clearly defined, typically 1 year. At each review:
			\begin{itemize}
				\item Verify the business need still exists
				\item Assess if compensating controls are effective
				\item Determine if the exception can be closed
				\item Re-approve if necessary
			\end{itemize}
		\end{block}
		\begin{alertblock}{Critical Trap}
			Exceptions are NOT permanent.
			Annual review = Minimum requirement.
		\end{alertblock}
	\end{frame}
	
	\begin{frame}{Q22: Vendor Offboarding Risk}
		\begin{block}{Question}
			What is the GREATEST risk during vendor offboarding?
		\end{block}
		\begin{enumerate}[(A)]
			\item Loss of service availability
			\item Data remanence (data not properly destroyed)
			\item Increased costs during transition
			\item Loss of institutional knowledge
		\end{enumerate}
	\end{frame}
	
	\begin{frame}{Q22: Explanation}
		\begin{block}{Correct Answer: B}
			\textbf{Data remanence} is the greatest risk because:
			\begin{itemize}
				\item Vendor may retain copies of your data
				\item Data may be in backups not properly deleted
				\item Vendor may reuse storage without sanitization
				\item Regulatory violations if data is not properly destroyed
				\item Reputational damage if data leaks post-termination
			\end{itemize}
		\end{block}
		\begin{alertblock}{Critical Trap}
			Offboarding = Data destruction is 1 priority.
			Certificate of destruction is required.
		\end{alertblock}
	\end{frame}
	
	% ============================================================
	% SECTION 4: CONTROL DESIGN AND IMPLEMENTATION - HARD QUESTIONS
	% ============================================================
	
	\section{4. Control Design and Implementation}
	
	\begin{frame}{Q23: Preventive vs. Deterrent}
		\begin{block}{Question}
			An organization installs signs warning of prosecution for unauthorized access. What type of control is this?
		\end{block}
		\begin{enumerate}[(A)]
			\item Preventive
			\item Detective
			\item Deterrent
			\item Corrective
		\end{enumerate}
	\end{frame}
	
	\begin{frame}{Q23: Explanation}
		\begin{block}{Correct Answer: C}
			\textbf{Deterrent controls} discourage potential attackers through psychological means. Signs warning of prosecution don't physically prevent access (preventive) but discourage it by threatening consequences.
		\end{block}
		\begin{alertblock}{Critical Trap}
			Deterrent = Discourages through fear/psychology.
			Preventive = Physically blocks through mechanics.
		\end{alertblock}
	\end{frame}
	
	\begin{frame}{Q24: Detective vs. Corrective Trap}
		\begin{block}{Question}
			An IDS detects an intrusion and automatically blocks the malicious IP. Is this detective or corrective?
		\end{block}
		\begin{enumerate}[(A)]
			\item Detective only
			\item Corrective only
			\item Both detective and corrective
			\item Preventive
		\end{enumerate}
	\end{frame}
	
	\begin{frame}{Q24: Explanation}
		\begin{block}{Correct Answer: C}
			The system performs \textbf{both} functions:
			\begin{itemize}
				\item \textbf{Detective:} It identifies the intrusion
				\item \textbf{Corrective:} It blocks the malicious IP (corrects the situation)
			\end{itemize}
			This is common in modern IDPS systems - they detect AND respond. The detection is detective; the blocking is corrective.
		\end{block}
		\begin{alertblock}{Critical Trap}
			Modern controls often combine multiple functions.
			IDS with auto-block = Detective + Corrective.
		\end{alertblock}
	\end{frame}
	
	\begin{frame}{Q25: Compensating Control Limitations}
		\begin{block}{Question}
			A compensating control is implemented because the primary control is not feasible. What must be TRUE for the compensating control to be acceptable?
		\end{block}
		\begin{enumerate}[(A)]
			\item It must be cheaper than the primary control
			\item It must provide equivalent or superior risk reduction
			\item It must be approved by the IT department
			\item It must be automated
		\end{enumerate}
	\end{frame}
	
	\begin{frame}{Q25: Explanation}
		\begin{block}{Correct Answer: B}
			A compensating control must provide \textbf{equivalent or superior risk reduction} to the primary control. Cost is not the deciding factor - effectiveness is. The control must address the same risk at an acceptable level.
		\end{block}
		\begin{alertblock}{Critical Trap}
			Compensating control = Same risk reduction, different method.
			Cheaper but less effective = NOT acceptable.
		\end{alertblock}
	\end{frame}
	
	\begin{frame}{Q26: Control Testing Frequency}
		\begin{block}{Question}
			How often should controls be tested?
		\end{block}
		\begin{enumerate}[(A)]
			\item Annually for all controls
			\item Based on risk criticality and regulatory requirements
			\item Only when there is a change
			\item Quarterly for all controls
		\end{enumerate}
	\end{frame}
	
	\begin{frame}{Q26: Explanation}
		\begin{block}{Correct Answer: B}
			Control testing frequency should be based on:
			\begin{itemize}
				\item \textbf{Risk criticality:} High-risk controls tested more frequently
				\item \textbf{Regulatory requirements:} Some regulations mandate specific frequencies
				\item \textbf{Control stability:} Stable controls may need less frequent testing
				\item \textbf{Change frequency:} Controls subject to change need more testing
			\end{itemize}
		\end{block}
		\begin{alertblock}{Critical Trap}
			One size does NOT fit all.
			Risk-based frequency > Fixed schedule.
		\end{alertblock}
	\end{frame}
	
	\begin{frame}{Q27: Control Implementation Failure}
		\begin{block}{Question}
			A control was implemented but is not operating as designed. This is an example of:
		\end{block}
		\begin{enumerate}[(A)]
			\item Control design deficiency
			\item Control operating deficiency
			\item Control gap
			\item Control redundancy
		\end{enumerate}
	\end{frame}
	
	\begin{frame}{Q27: Explanation}
		\begin{block}{Correct Answer: B}
			\textbf{Operating deficiency} = Control is designed correctly but not operating effectively. \textbf{Design deficiency} = Control is not designed to address the risk. The distinction is critical for remediation.
		\end{block}
		\begin{alertblock}{Critical Trap}
			Design deficiency = Wrong control for the risk.
			Operating deficiency = Right control, not working.
		\end{alertblock}
	\end{frame}
	
	\begin{frame}{Q28: Parallel Changeover Risk}
		\begin{block}{Question}
			What is the PRIMARY risk of parallel changeover?
		\end{block}
		\begin{enumerate}[(A)]
			\item Data inconsistency between systems
			\item Complete system failure
			\item User resistance to change
			\item Extended implementation timeline
		\end{enumerate}
	\end{frame}
	
	\begin{frame}{Q28: Explanation}
		\begin{block}{Correct Answer: A}
			\textbf{Data inconsistency} is the primary risk because:
			\begin{itemize}
				\item Both systems are processing transactions
				\item Data may be entered in one but not the other
				\item Synchronization failures can occur
				\item Reconciliation becomes complex
			\end{itemize}
			While parallel is the safest for availability, data consistency is the trade-off.
		\end{block}
		\begin{alertblock}{Critical Trap}
			Parallel = Safest for availability, riskiest for data consistency.
			Data reconciliation is critical.
		\end{alertblock}
	\end{frame}
	
	\begin{frame}{Q29: Abrupt Changeover Appropriate When}
		\begin{block}{Question}
			When is abrupt changeover acceptable?
		\end{block}
		\begin{enumerate}[(A)]
			\item When the system is critical to business operations
			\item When the change is minor and rollback is easy
			\item When users are well-trained
			\item When the vendor recommends it
		\end{enumerate}
	\end{frame}
	
	\begin{frame}{Q29: Explanation}
		\begin{block}{Correct Answer: B}
			Abrupt changeover is only acceptable when:
			\begin{itemize}
				\item The change is \textbf{minor} (low risk)
				\item \textbf{Rollback is easy} (can revert quickly)
				\item \textbf{Business impact is minimal}
				\item No critical business processes are affected
			\end{itemize}
		\end{block}
		\begin{alertblock}{Critical Trap}
			Abrupt = Fastest but highest risk.
			Only for non-critical, easily reversible changes.
		\end{alertblock}
	\end{frame}
	
	\begin{frame}{Q30: Post-Implementation Review Timing}
		\begin{block}{Question}
			When should a post-implementation review occur?
		\end{block}
		\begin{enumerate}[(A)]
			\item Immediately after go-live
			\item After the system has been stable for a period
			\item Only if there were issues
			\item Before the next audit
		\end{enumerate}
	\end{frame}
	
	\begin{frame}{Q30: Explanation}
		\begin{block}{Correct Answer: B}
			The post-implementation review should occur \textbf{after the system has been stable for a period} (typically 3-6 months) to:
			\begin{itemize}
				\item Allow issues to surface
				\item Gather meaningful user feedback
				\item Assess actual performance vs. expected
				\item Identify lessons learned
			\end{itemize}
		\end{block}
		\begin{alertblock}{Critical Trap}
			Too early = Issues haven't surfaced.
			Too late = Lessons forgotten.
			Wait for stability, then review.
		\end{alertblock}
	\end{frame}
	
	% ============================================================
	% SECTION 5: MONITORING AND REPORTING - HARD QUESTIONS
	% ============================================================
	
	\section{5. Monitoring and Reporting}
	
	\begin{frame}{Q31: KRI vs. KPI vs. KCI}
		\begin{block}{Question}
			The percentage of employees who completed security awareness training is a:
		\end{block}
		\begin{enumerate}[(A)]
			\item Key Risk Indicator (KRI)
			\item Key Performance Indicator (KPI)
			\item Key Control Indicator (KCI)
			\item None of the above
		\end{enumerate}
	\end{frame}
	
	\begin{frame}{Q31: Explanation}
		\begin{block}{Correct Answer: B}
			Training completion is a \textbf{KPI} because:
			\begin{itemize}
				\item It measures \textbf{performance} of the training program
				\item It looks \textbf{backward} at what has been accomplished
				\item It doesn't predict risk (that would be KRI)
				\item It doesn't measure control effectiveness directly (that would be KCI)
			\end{itemize}
		\end{block}
		\begin{alertblock}{Critical Trap}
			Training completion = KPI (Performance).
			Phishing click rate = KRI (Risk predictor).
			Access review completion = KCI (Control effectiveness).
		\end{alertblock}
	\end{frame}
	
	\begin{frame}{Q32: KRI Threshold Breach}
		\begin{block}{Question}
			A KRI has a threshold of 10%. Current value is 12%. What should happen FIRST?
		\end{block}
		\begin{enumerate}[(A)]
			\item Escalate to senior management
			\item Investigate the cause of the breach
			\item Implement additional controls
			\item Update the risk register
		\end{enumerate}
	\end{frame}
	
	\begin{frame}{Q32: Explanation}
		\begin{block}{Correct Answer: B}
			\textbf{Investigate first} before taking action. The breach could be:
			\begin{itemize}
				\item A data quality issue (false alarm)
				\item A temporary spike (not a trend)
				\item A genuine risk increase requiring response
			\end{itemize}
			Understanding the cause determines the appropriate response.
		\end{block}
		\begin{alertblock}{Critical Trap}
			Threshold breach not equal to Immediate escalation.
			Investigate first, then respond appropriately.
		\end{alertblock}
	\end{frame}
	
	\begin{frame}{Q33: Dashboard Audience}
		\begin{block}{Question}
			Which report format is MOST appropriate for the Board of Directors?
		\end{block}
		\begin{enumerate}[(A)]
			\item Detailed technical vulnerability report
			\item One-page executive summary with key metrics
			\item Full risk register with all 500 risks
			\item Raw SIEM logs
		\end{enumerate}
	\end{frame}
	
	\begin{frame}{Q33: Explanation}
		\begin{block}{Correct Answer: B}
			The Board needs \textbf{concise, business-focused information}:
			\begin{itemize}
				\item One page maximum
				\item Key metrics (KRIs/KPIs)
				\item Business impact, not technical details
				\item Trends and actionable insights
				\item Comparison to risk appetite
			\end{itemize}
		\end{block}
		\begin{alertblock}{Critical Trap}
			Board not equal to Technical audience.
			Business language, not technical jargon.
		\end{alertblock}
	\end{frame}
	
	\begin{frame}{Q34: SMART KRI Criteria}
		\begin{block}{Question}
			A KRI is defined as "number of security incidents." What is the primary weakness of this KRI?
		\end{block}
		\begin{enumerate}[(A)]
			\item It is not measurable
			\item It lacks specificity and thresholds
			\item It is too difficult to track
			\item It is not relevant to security
		\end{enumerate}
	\end{frame}
	
	\begin{frame}{Q34: Explanation}
		\begin{block}{Correct Answer: B}
			"Number of security incidents" lacks:
			\begin{itemize}
				\item \textbf{Specificity:} What type of incidents? What severity?
				\item \textbf{Thresholds:} What number triggers action?
				\item \textbf{Timeframe:} Per day? Per month? Trend?
				\item \textbf{Actionability:} What response is expected?
			\end{itemize}
			A good KRI: "Number of critical severity incidents per month exceeding threshold of 3."
		\end{block}
		\begin{alertblock}{Critical Trap}
			Vague KRI = Useless KRI.
			Specific + Measurable + Threshold + Actionable.
		\end{alertblock}
	\end{frame}
	
	\begin{frame}{Q35: Control Monitoring Purpose}
		\begin{block}{Question}
			The PRIMARY purpose of control monitoring is to:
		\end{block}
		\begin{enumerate}[(A)]
			\item Verify the control is working as designed
			\item Verify the control is effectively reducing risk
			\item Verify the control is cost-effective
			\item Verify the control is automated
		\end{enumerate}
	\end{frame}
	
	\begin{frame}{Q35: Explanation}
		\begin{block}{Correct Answer: B}
			The primary purpose is to verify \textbf{risk reduction effectiveness}, not just operation. A control can be "working as designed" but not actually reducing risk if:
			\begin{itemize}
				\item The design is flawed
				\item The threat landscape has changed
				\item The control is redundant
			\end{itemize}
		\end{block}
		\begin{alertblock}{Critical Trap}
			Control working not equal to Control effective.
			Monitor effectiveness, not just operation.
		\end{alertblock}
	\end{frame}
	
	\begin{frame}{Q36: Heat Map Limitation}
		\begin{block}{Question}
			What is the PRIMARY limitation of a risk heat map?
		\end{block}
		\begin{enumerate}[(A)]
			\item It is too complex for executives
			\item It is qualitative and subjective
			\item It cannot show multiple risks
			\item It requires too much data
		\end{enumerate}
	\end{frame}
	
	\begin{frame}{Q36: Explanation}
		\begin{block}{Correct Answer: B}
			Heat maps are \textbf{qualitative and subjective}:
			\begin{itemize}
				\item "High" impact means different things to different people
				\item No quantification of actual loss
				\item Difficult to compare across business units
				\item Can oversimplify complex risks
			\end{itemize}
			Dashboards with quantitative metrics address this limitation.
		\end{block}
		\begin{alertblock}{Critical Trap}
			Heat map = Qualitative = Subjective.
			Dashboard = Quantitative + Qualitative = More objective.
		\end{alertblock}
	\end{frame}
	
	\begin{frame}{Q37: Reporting Frequency}
		\begin{block}{Question}
			How often should risk reports be provided to senior management?
		\end{block}
		\begin{enumerate}[(A)]
			\item Annually
			\item Quarterly
			\item Monthly
			\item Based on risk velocity and organizational needs
		\end{enumerate}
	\end{frame}
	
	\begin{frame}{Q37: Explanation}
		\begin{block}{Correct Answer: D}
			Reporting frequency should be \textbf{risk-based}:
			\begin{itemize}
				\item \textbf{High-velocity risks:} More frequent (monthly/weekly)
				\item \textbf{Stable risks:} Less frequent (quarterly)
				\item \textbf{Regulatory requirements:} May mandate frequency
				\item \textbf{Organizational culture:} Some prefer more frequent updates
			\end{itemize}
		\end{block}
		\begin{alertblock}{Critical Trap}
			One size does NOT fit all.
			Risk velocity determines reporting frequency.
		\end{alertblock}
	\end{frame}
	
	\begin{frame}{Q38: Exception Reporting}
		\begin{block}{Question}
			What should be included in an exception report to management?
		\end{block}
		\begin{enumerate}[(A)]
			\item Only the list of exceptions
			\item Exceptions, compensating controls, risk owner, and expiration date
			\item Only exceptions that exceed 1 year
			\item Only exceptions approved by IT
		\end{enumerate}
	\end{frame}
	
	\begin{frame}{Q38: Explanation}
		\begin{block}{Correct Answer: B}
			Exception reports should include:
			\begin{itemize}
				\item \textbf{Description:} What is the exception?
				\item \textbf{Compensating controls:} How is risk being managed?
				\item \textbf{Risk owner:} Who accepted the risk?
				\item \textbf{Expiration date:} When must it be renewed?
				\item \textbf{Business justification:} Why is it needed?
			\end{itemize}
		\end{block}
		\begin{alertblock}{Critical Trap}
			Exception report not equal to Simple list.
			Must include risk context and management details.
		\end{alertblock}
	\end{frame}
	
	\begin{frame}{Q39: KCI vs. KRI}
		\begin{block}{Question}
			The number of access reviews completed on time is a:
		\end{block}
		\begin{enumerate}[(A)]
			\item Key Risk Indicator (KRI)
			\item Key Performance Indicator (KPI)
			\item Key Control Indicator (KCI)
			\item None of the above
		\end{enumerate}
	\end{frame}
	
	\begin{frame}{Q39: Explanation}
		\begin{block}{Correct Answer: C}
			Access reviews are a \textbf{control}. The percentage completed on time measures \textbf{control effectiveness} = KCI.
			\begin{itemize}
				\item KRI: Predicts risk (e.g., number of orphan accounts)
				\item KPI: Measures performance (e.g., time to complete reviews)
				\item KCI: Measures control effectiveness (e.g., % reviews on time)
			\end{itemize}
		\end{block}
		\begin{alertblock}{Critical Trap}
			KCI = Control effectiveness.
			If it measures a control, it's a KCI.
		\end{alertblock}
	\end{frame}
	
	\begin{frame}{Q40: Log Retention}
		\begin{block}{Question}
			How long should security logs be retained?
		\end{block}
		\begin{enumerate}[(A)]
			\item 30 days
			\item 90 days
			\item Based on regulatory requirements and forensic needs
			\item Indefinitely
		\end{enumerate}
	\end{frame}
	
	\begin{frame}{Q40: Explanation}
		\begin{block}{Correct Answer: C}
			Log retention should be based on:
			\begin{itemize}
				\item \textbf{Regulatory requirements:} Some industries mandate 1-7 years
				\item \textbf{Forensic needs:} Enough time to investigate incidents
				\item \textbf{Storage costs:} Balance cost vs. benefit
				\item \textbf{Legal hold requirements:} Litigation may require preservation
			\end{itemize}
		\end{block}
		\begin{alertblock}{Critical Trap}
			No universal log retention period.
			Regulatory + Forensic + Legal = Retention period.
		\end{alertblock}
	\end{frame}
	
	% ============================================================
	% SECTION 6: ADVANCED TRAPS AND SCENARIOS
	% ============================================================
	
	\section{6. Advanced Traps and Scenarios}
	
	\begin{frame}{Q41: Risk Response Selection Trap}
		\begin{block}{Question}
			A risk has been assessed as "High" (likelihood: 80%, impact: \$500K). The cost to mitigate is \$100K. The company has a risk appetite of \$200K. What is the BEST response?
		\end{block}
		\begin{enumerate}[(A)]
			\item Accept the risk because mitigation cost is high
			\item Mitigate because cost (\$100K) < risk exposure (\$400K)
			\item Transfer the risk through insurance
			\item Avoid the risk entirely
		\end{enumerate}
	\end{frame}
	
	\begin{frame}{Q41: Explanation}
		\begin{block}{Correct Answer: B}
			\textbf{Mitigation is justified:}
			\begin{itemize}
				\item Expected loss = 80% × \$500K = \$400K
				\item Control cost = \$100K
				\item Net benefit = \$300K
				\item Residual risk would be below appetite
			\end{itemize}
		\end{block}
		\begin{alertblock}{Critical Trap}
			Expected loss = Likelihood × Impact.
			Mitigation justified when control cost < expected loss.
		\end{alertblock}
	\end{frame}
	
	\begin{frame}{Q42: Risk Owner Authority Trap}
		\begin{block}{Question}
			A risk owner wants to accept a risk that exceeds the organization's risk appetite. They argue that their business unit has special circumstances. What should happen?
		\end{block}
		\begin{enumerate}[(A)]
			\item Allow the acceptance because they are the risk owner
			\item Escalate to senior management for approval
			\item Automatically transfer the risk
			\item Automatically avoid the risk
		\end{enumerate}
	\end{frame}
	
	\begin{frame}{Q42: Explanation}
		\begin{block}{Correct Answer: B}
			Risk acceptance above the \textbf{enterprise risk appetite} requires \textbf{executive approval}. A business unit cannot unilaterally accept risk that exceeds organizational tolerance levels. This ensures:
			\begin{itemize}
				\item Enterprise-level perspective on risk
				\item Consistency in risk-taking across BUs
				\item Board oversight of significant risk
			\end{itemize}
		\end{block}
		\begin{alertblock}{Critical Trap}
			Risk owner not equal to Unlimited authority.
			Exceeding risk appetite = Escalation required.
		\end{alertblock}
	\end{frame}
	
	\begin{frame}{Q43: Third-Party Assessment Trap}
		\begin{block}{Question}
			A vendor provides a SOC 2 Type II report with no exceptions. Can you rely on this for vendor security assessment?
		\end{block}
		\begin{enumerate}[(A)]
			\item Yes, SOC 2 Type II with no exceptions means the vendor is secure
			\item Yes, if the scope matches your requirements
			\item No, you must always conduct your own assessment
			\item No, SOC 2 is not relevant for vendor assessment
		\end{enumerate}
	\end{frame}
	
	\begin{frame}{Q43: Explanation}
		\begin{block}{Correct Answer: B}
			\textbf{It depends on scope.} A SOC 2 Type II report is valuable IF:
			\begin{itemize}
				\item The scope covers the services you're using
				\item The Trust Services Criteria are relevant
				\item The period is recent (within 12 months)
				\item No exceptions or compensating controls are acceptable
			\end{itemize}
			You may still need supplemental assessment for gaps.
		\end{block}
		\begin{alertblock}{Critical Trap}
			SOC 2 = Useful, but check scope first.
			Wrong scope = Wrong assurance.
		\end{alertblock}
	\end{frame}
	
	\begin{frame}{Q44: Control Monitoring Trap}
		\begin{block}{Question}
			A control has been operating effectively for 3 years. Can the frequency of testing be reduced?
		\end{block}
		\begin{enumerate}[(A)]
			\item No, controls must always be tested annually
			\item Yes, if the control environment is stable and risk is low
			\item No, regulatory requirements mandate annual testing
			\item Yes, after 3 years of effective operation, testing can stop
		\end{enumerate}
	\end{frame}
	
	\begin{frame}{Q44: Explanation}
		\begin{block}{Correct Answer: B}
			Testing frequency can be \textbf{reduced} if:
			\begin{itemize}
				\item Control has a strong track record (3 years effective)
				\item The control environment is stable
				\item Risk is low/stable
				\item No significant changes to the process
				\item Regulatory requirements allow it
			\end{itemize}
			However, monitoring should never stop completely.
		\end{block}
		\begin{alertblock}{Critical Trap}
			Strong track record = Reduced testing, NOT no testing.
			Risk-based frequency adjustments are acceptable.
		\end{alertblock}
	\end{frame}
	
	\begin{frame}{Q45: Risk Acceptance Trap}
		\begin{block}{Question}
			A CEO verbally accepts a high risk without documentation. Is this valid risk acceptance?
		\end{block}
		\begin{enumerate}[(A)]
			\item Yes, CEO has ultimate authority
			\item No, risk acceptance must be documented
			\item Yes, if the CEO communicates it to the board
			\item No, only the board can accept risk
		\end{enumerate}
	\end{frame}
	
	\begin{frame}{Q45: Explanation}
		\begin{block}{Correct Answer: B}
			Risk acceptance must be \textbf{formally documented} and include:
			\begin{itemize}
				\item Risk description and impact
				\item Rationale for acceptance
				\item Risk owner signature
				\item Duration of acceptance
				\item Review schedule
			\end{itemize}
			Verbal acceptance provides no audit trail and is not defensible.
		\end{block}
		\begin{alertblock}{Critical Trap}
			Verbal acceptance = No accountability.
			Documentation is mandatory.
		\end{alertblock}
	\end{frame}
	
	\begin{frame}{Q46: Control Implementation Trap}
		\begin{block}{Question}
			A control has been designed and implemented. What must happen before it can be considered effective?
		\end{block}
		\begin{enumerate}[(A)]
			\item The vendor must certify it
			\item It must be tested and verified
			\item It must be approved by the board
			\item It must be automated
		\end{enumerate}
	\end{frame}
	
	\begin{frame}{Q46: Explanation}
		\begin{block}{Correct Answer: B}
			A control is not effective until it has been \textbf{tested and verified}. Implementation not equal to Effectiveness. Testing confirms:
			\begin{itemize}
				\item The control operates as designed
				\item The control actually reduces risk
				\item There are no gaps or weaknesses
				\item The control is sustainable
			\end{itemize}
		\end{block}
		\begin{alertblock}{Critical Trap}
			Implementation not equal to Effectiveness.
			Testing is required to confirm effectiveness.
		\end{alertblock}
	\end{frame}
	
	\begin{frame}{Q47: KRI Maintenance Trap}
		\begin{block}{Question}
			Why must KRIs be maintained and updated regularly?
		\end{block}
		\begin{enumerate}[(A)]
			\item To keep them complex and sophisticated
			\item Because threats and vulnerabilities change over time
			\item To justify the risk management budget
			\item To satisfy audit requirements
		\end{enumerate}
	\end{frame}
	
	\begin{frame}{Q47: Explanation}
		\begin{block}{Correct Answer: B}
			KRIs must evolve because the \textbf{risk environment changes}:
			\begin{itemize}
				\item New threats emerge
				\item Vulnerabilities are discovered/patch
				\item Business objectives change
				\item Regulatory requirements evolve
				\item Technology landscape shifts
			\end{itemize}
			Stale KRIs provide false assurance.
		\end{block}
		\begin{alertblock}{Critical Trap}
			KRIs are NOT set-and-forget.
			Regular review and update is required.
		\end{alertblock}
	\end{frame}
	
	\begin{frame}{Q48: Risk Reporting Trap}
		\begin{block}{Question}
			A risk report shows all risks are "green" (low). What should the risk manager do?
		\end{block}
		\begin{enumerate}[(A)]
			\item Celebrate - the organization is secure
			\item Investigate whether the assessment is accurate or overly optimistic
			\item Reduce the risk management budget
			\item Focus on other priorities
		\end{enumerate}
	\end{frame}
	
	\begin{frame}{Q48: Explanation}
		\begin{block}{Correct Answer: B}
			All "green" is a \textbf{red flag}. It could indicate:
			\begin{itemize}
				\item Risk assessment is too lenient
				\item Risks are not being identified
				\item Thresholds are set too high
				\item The organization is in denial
				\item Complacency has set in
			\end{itemize}
			A healthy risk portfolio usually has some amber/red items.
		\end{block}
		\begin{alertblock}{Critical Trap}
			All green = Suspicious.
			Investigate, don't celebrate.
		\end{alertblock}
	\end{frame}
	
	\begin{frame}{Q49: Control Testing Trap}
		\begin{block}{Question}
			A control test reveals one exception out of 100 samples. What should the risk manager conclude?
		\end{block}
		\begin{enumerate}[(A)]
			\item The control is effective (99% success rate)
			\item The control has a deficiency that requires investigation
			\item The exception should be ignored (statistical anomaly)
			\item The control should be immediately replaced
		\end{enumerate}
	\end{frame}
	
	\begin{frame}{Q49: Explanation}
		\begin{block}{Correct Answer: B}
			Any exception requires \textbf{investigation}:
			\begin{itemize}
				\item What caused the exception?
				\item Is it a design flaw or operating issue?
				\item Is it an isolated event or systemic?
				\item Are compensating controls in place?
			\end{itemize}
			Even 1% failure rate can have significant impact for critical controls.
		\end{block}
		\begin{alertblock}{Critical Trap}
			Any exception = Investigate.
			Don't dismiss as statistical anomaly.
		\end{alertblock}
	\end{frame}
	
	\begin{frame}{Q50: Risk Treatment Trap}
		\begin{block}{Question}
			A risk has been mitigated to within risk appetite. What is the NEXT step?
		\end{block}
		\begin{enumerate}[(A)]
			\item Close the risk and stop monitoring
			\item Continue monitoring to ensure risk remains within appetite
			\item Transfer the remaining risk
			\item Accept the residual risk
		\end{enumerate}
	\end{frame}
	
	\begin{frame}{Q50: Explanation}
		\begin{block}{Correct Answer: B}
			After mitigation, \textbf{continuous monitoring} is required because:
			\begin{itemize}
				\item Risk levels can change over time
				\item Controls can degrade or fail
				\item New threats may emerge
				\item Business changes may alter risk profile
			\end{itemize}
			Residual risk is not static.
		\end{block}
		\begin{alertblock}{Critical Trap}
			Mitigated not equal to Forgotten.
			Continuous monitoring is essential.
		\end{alertblock}
	\end{frame}
	
	% ============================================================
	% SECTION 7: MORE ADVANCED QUESTIONS
	% ============================================================
	
	\section{7. More Advanced Questions}
	
	\begin{frame}{Q51: Risk Appetite vs. Tolerance Trap}
		\begin{block}{Question}
			The risk appetite is \$100K. The risk tolerance is \$150K. A risk is currently at \$120K. What does this mean?
		\end{block}
		\begin{enumerate}[(A)]
			\item The risk is within appetite
			\item The risk exceeds appetite but is within tolerance
			\item The risk exceeds both appetite and tolerance
			\item The risk is within tolerance but exceeds capacity
		\end{enumerate}
	\end{frame}
	
	\begin{frame}{Q51: Explanation}
		\begin{block}{Correct Answer: B}
			\textbf{Risk appetite} = \$100K (desired level)\\
			\textbf{Risk tolerance} = \$150K (acceptable deviation)\\
			\textbf{Current risk} = \$120K
			
			The risk is:
			\begin{itemize}
				\item Above appetite (\$120K > \$100K) - requires attention
				\item Within tolerance (\$120K < \$150K) - acceptable with monitoring
			\end{itemize}
		\end{block}
		\begin{alertblock}{Critical Trap}
			Appetite = Target.
			Tolerance = Acceptable range.
			Between them = Monitor closely.
		\end{alertblock}
	\end{frame}
	
	\begin{frame}{Q52: Risk Capacity Trap}
		\begin{block}{Question}
			The risk capacity is \$500K. A risk is at \$450K. What is the recommended action?
		\end{block}
		\begin{enumerate}[(A)]
			\item No action needed - within capacity
			\item Immediate action to reduce risk below appetite
			\item Escalate to senior management and reduce risk urgently
			\item Transfer the risk immediately
		\end{enumerate}
	\end{frame}
	
	\begin{frame}{Q52: Explanation}
		\begin{block}{Correct Answer: C}
			Approaching \textbf{risk capacity} (\$450K vs. \$500K) is a \textbf{critical situation}:
			\begin{itemize}
				\item Capacity = maximum risk before existence is threatened
				\item Being at 90% of capacity is extremely dangerous
				\item Any additional risk could threaten the organization
				\item Immediate escalation and action required
			\end{itemize}
		\end{block}
		\begin{alertblock}{Critical Trap}
			Approaching capacity = Emergency.
			Capacity breach = Existential threat.
		\end{alertblock}
	\end{frame}
	
	\begin{frame}{Q53: Inherent vs. Residual Trap}
		\begin{block}{Question}
			A risk has an inherent risk of \$1M. Controls costing \$200K reduce this to \$300K. What is the residual risk?
		\end{block}
		\begin{enumerate}[(A)]
			\item \$1M
			\item \$300K
			\item \$500K
			\item \$800K
		\end{enumerate}
	\end{frame}
	
	\begin{frame}{Q53: Explanation}
		\begin{block}{Correct Answer: B}
			\textbf{Residual risk = \$300K} (the risk remaining after controls). The calculation:
			\begin{itemize}
				\item Inherent risk = \$1M
				\item Control cost = \$200K (not subtracted from risk)
				\item Residual risk = \$300K (risk after controls)
			\end{itemize}
			Control cost is a separate consideration (cost-benefit analysis), not a risk reduction.
		\end{block}
		\begin{alertblock}{Critical Trap}
			Residual risk = Risk after controls, NOT risk minus control cost.
		\end{alertblock}
	\end{frame}
	
	\begin{frame}{Q54: Control Cost Trap}
		\begin{block}{Question}
			A control costs \$100K and reduces risk from \$500K to \$100K. What is the ROI of the control?
		\end{block}
		\begin{enumerate}[(A)]
			\item 100%
			\item 300%
			\item 400%
			\item Cannot be determined
		\end{enumerate}
	\end{frame}
	
	\begin{frame}{Q54: Explanation}
		\begin{block}{Correct Answer: B}
			\textbf{ROI = 300\%}
			\begin{itemize}
				\item Risk reduction = \$500K - \$100K = \$400K
				\item Control cost = \$100K
				\item Net benefit = \$400K - \$100K = \$300K
				\item ROI = Net benefit / Control cost = \$300K / \$100K = 300\%
			\end{itemize}
		\end{block}
		\begin{alertblock}{Critical Trap}
			ROI = (Risk Reduction - Control Cost) / Control Cost.
			Not just risk reduction / control cost.
		\end{alertblock}
	\end{frame}
	
	\begin{frame}{Q55: Detection vs. Prevention Trap}
		\begin{block}{Question}
			An organization installs motion detectors. What type of control is this?
		\end{block}
		\begin{enumerate}[(A)]
			\item Preventive
			\item Detective
			\item Deterrent
			\item Corrective
		\end{enumerate}
	\end{frame}
	
	\begin{frame}{Q55: Explanation}
		\begin{block}{Correct Answer: B}
			Motion detectors are \textbf{detective} controls:
			\begin{itemize}
				\item They detect movement (intrusion)
				\item They do NOT prevent entry
				\item They alert to the presence of an intruder
				\item Detection triggers response (corrective)
			\end{itemize}
			They may have a deterrent effect, but their primary function is detection.
		\end{block}
		\begin{alertblock}{Critical Trap}
			Motion detector = Detective.
			It detects, it doesn't prevent.
		\end{alertblock}
	\end{frame}
	
	\begin{frame}{Q56: Risk Acceptance Authority Trap}
		\begin{block}{Question}
			Who has the authority to accept a risk that exceeds the enterprise risk appetite?
		\end{block}
		\begin{enumerate}[(A)]
			\item Risk owner
			\item Control owner
			\item Board of directors or equivalent
			\item IT manager
		\end{enumerate}
	\end{frame}
	
	\begin{frame}{Q56: Explanation}
		\begin{block}{Correct Answer: C}
			Risk acceptance \textbf{above enterprise risk appetite} requires \textbf{board-level approval} because:
			\begin{itemize}
				\item It affects the entire organization
				\item It may impact shareholder value
				\item It requires strategic perspective
				\item It may have regulatory implications
			\end{itemize}
		\end{block}
		\begin{alertblock}{Critical Trap}
			Board sets risk appetite.
			Board approves exceptions to appetite.
		\end{alertblock}
	\end{frame}
	
	\begin{frame}{Q57: Control Design Trap}
		\begin{block}{Question}
			A control is designed to prevent unauthorized access. However, it is too complex for users, so they bypass it. What type of deficiency is this?
		\end{block}
		\begin{enumerate}[(A)]
			\item Operating deficiency
			\item Design deficiency
			\item Both design and operating deficiency
			\item No deficiency
		\end{enumerate}
	\end{frame}
	
	\begin{frame}{Q57: Explanation}
		\begin{block}{Correct Answer: B}
			This is a \textbf{design deficiency} because:
			\begin{itemize}
				\item The control was designed poorly (too complex)
				\item User bypass is a predictable consequence
				\item The design failed to consider usability
				\item A better design would balance security and usability
			\end{itemize}
			The control doesn't work because of how it was designed.
		\end{block}
		\begin{alertblock}{Critical Trap}
			Design flaw causes operating failure.
			Root cause = Design deficiency.
		\end{alertblock}
	\end{frame}
	
	\begin{frame}{Q58: Risk Register Trap}
		\begin{block}{Question}
			What is the MOST important element of a risk register entry?
		\end{block}
		\begin{enumerate}[(A)]
			\item Risk description
			\item Risk owner
			\item Risk score
			\item Mitigation plan
		\end{enumerate}
	\end{frame}
	
	\begin{frame}{Q58: Explanation}
		\begin{block}{Correct Answer: B}
			\textbf{Risk owner} is most important because:
			\begin{itemize}
				\item Without an owner, no one is accountable
				\item Other elements (description, score, plan) are useless without ownership
				\item Owner drives the risk management process
				\item Owner ensures the risk is monitored and addressed
			\end{itemize}
		\end{block}
		\begin{alertblock}{Critical Trap}
			No owner = No accountability = No risk management.
			Owner is the foundation.
		\end{alertblock}
	\end{frame}
	
	\begin{frame}{Q59: Third-Party Risk Trap}
		\begin{block}{Question}
			A vendor has a SOC 2 Type II report with no exceptions. Your company uses the vendor for a service NOT covered in the SOC 2 scope. What should you do?
		\end{block}
		\begin{enumerate}[(A)]
			\item Rely on the SOC 2 report anyway
			\item Conduct additional assessment for the out-of-scope service
			\item Reject the vendor entirely
			\item Request a new SOC 2 report
		\end{enumerate}
	\end{frame}
	
	\begin{frame}{Q59: Explanation}
		\begin{block}{Correct Answer: B}
			The SOC 2 report only covers the \textbf{scope defined in the report}. If your service is out of scope, you must:
			\begin{itemize}
				\item Conduct your own assessment of that service
				\item Request a scope expansion from the vendor
				\item Implement additional controls
				\item Document the gap and residual risk
			\end{itemize}
		\end{block}
		\begin{alertblock}{Critical Trap}
			SOC 2 scope not equal to All services.
			Out of scope = Additional assessment required.
		\end{alertblock}
	\end{frame}
	
	\begin{frame}{Q60: Risk Response Documentation Trap}
		\begin{block}{Question}
			Why must risk responses be documented?
		\end{block}
		\begin{enumerate}[(A)]
			\item To satisfy auditors
			\item To provide evidence of due care and accountability
			\item To justify budget requests
			\item To track project progress
		\end{enumerate}
	\end{frame}
	
	\begin{frame}{Q60: Explanation}
		\begin{block}{Correct Answer: B}
			Documentation provides \textbf{evidence of due care}:
			\begin{itemize}
				\item Shows the organization took risk seriously
				\item Demonstrates accountability
				\item Protects against negligence claims
				\item Provides audit trail for decisions
				\item Enables continuous improvement
			\end{itemize}
		\end{block}
		\begin{alertblock}{Critical Trap}
			Undocumented = Undefensible.
			Documentation = Due care evidence.
		\end{alertblock}
	\end{frame}
	
	% ============================================================
	% SECTION 8: FINAL ADVANCED QUESTIONS
	% ============================================================
	
	\section{8. Final Advanced Questions}
	
	\begin{frame}{Q61: Risk Velocity Trap}
		\begin{block}{Question}
			What is risk velocity and why does it matter?
		\end{block}
		\begin{enumerate}[(A)]
			\item Speed at which risk is identified
			\item Speed at which risk materializes
			\item Speed at which controls are implemented
			\item Speed at which risk is reported
		\end{enumerate}
	\end{frame}
	
	\begin{frame}{Q61: Explanation}
		\begin{block}{Correct Answer: B}
			\textbf{Risk velocity} = speed at which risk materializes. It matters because:
			\begin{itemize}
				\item High-velocity risks require faster response
				\item Affects monitoring frequency (higher velocity = more frequent)
				\item Impacts risk appetite (high velocity = lower tolerance)
				\item Determines urgency of mitigation
			\end{itemize}
		\end{block}
		\begin{alertblock}{Critical Trap}
			High velocity = Less time to react.
			Velocity determines response urgency.
		\end{alertblock}
	\end{frame}
	
	\begin{frame}{Q62: Key Risk Indicator Trap}
		\begin{block}{Question}
			A KRI is showing a consistent upward trend but is still below threshold. What should the risk manager do?
		\end{block}
		\begin{enumerate}[(A)]
			\item No action needed - below threshold
			\item Investigate the trend and prepare for potential threshold breach
			\item Wait for the threshold to be breached
			\item Report to the board immediately
		\end{enumerate}
	\end{frame}
	
	\begin{frame}{Q62: Explanation}
		\begin{block}{Correct Answer: B}
			\textbf{Trend analysis} is critical:
			\begin{itemize}
				\item Upward trend indicates increasing risk
				\item Threshold may be breached soon
				\item Preventive action can be taken now
				\item Understanding the cause enables proactive response
			\end{itemize}
		\end{block}
		\begin{alertblock}{Critical Trap}
			Below threshold not equal to No action.
			Trend is as important as current value.
		\end{alertblock}
	\end{frame}
	
	\begin{frame}{Q63: Control Monitoring Trap}
		\begin{block}{Question}
			A control has been automated. Does this eliminate the need for monitoring?
		\end{block}
		\begin{enumerate}[(A)]
			\item Yes, automated controls are self-monitoring
			\item No, automated controls can fail and must be monitored
			\item Yes, automation eliminates human error
			\item No, but monitoring frequency can be reduced
		\end{enumerate}
	\end{frame}
	
	\begin{frame}{Q63: Explanation}
		\begin{block}{Correct Answer: B}
			Automated controls \textbf{still require monitoring}:
			\begin{itemize}
				\item Automation can fail (bugs, configuration errors)
				\item Automated controls can be bypassed
				\item The automation logic may become outdated
				\item Regular testing confirms continued effectiveness
			\end{itemize}
		\end{block}
		\begin{alertblock}{Critical Trap}
			Automated not equal Self-monitoring.
			Automation reduces human error but doesn't eliminate need for oversight.
		\end{alertblock}
	\end{frame}
	
	\begin{frame}{Q64: Risk Acceptance Documentation Trap}
		\begin{block}{Question}
			A risk acceptance was documented but not periodically reviewed. What is the consequence?
		\end{block}
		\begin{enumerate}[(A)]
			\item No consequence - documentation exists
			\item The risk may have changed; old acceptance may not be valid
			\item The risk is automatically transferred
			\item The risk is automatically avoided
		\end{enumerate}
	\end{frame}
	
	\begin{frame}{Q64: Explanation}
		\begin{block}{Correct Answer: B}
			Without periodic review:
			\begin{itemize}
				\item The risk may have increased
				\item The business context may have changed
				\item New controls may make acceptance unnecessary
				\item The acceptance may have expired
			\end{itemize}
			Regular review ensures the acceptance remains valid and appropriate.
		\end{block}
		\begin{alertblock}{Critical Trap}
			Old acceptance = Potentially invalid.
			Annual review is mandatory.
		\end{alertblock}
	\end{frame}
	
	\begin{frame}{Q65: Control Testing Approach Trap}
		\begin{block}{Question}
			When testing a control, what is the BEST approach?
		\end{block}
		\begin{enumerate}[(A)]
			\item Test every transaction
			\item Use statistical sampling and document methodology
			\item Test only the most recent transactions
			\item Rely on management's assurance
		\end{enumerate}
	\end{frame}
	
	\begin{frame}{Q65: Explanation}
		\begin{block}{Correct Answer: B}
			\textbf{Statistical sampling} is best because:
			\begin{itemize}
				\item Testing every transaction is impractical
				\item Proper sampling provides reasonable assurance
				\item Documented methodology enables repeatability
				\item Results can be defended to auditors
			\end{itemize}
		\end{block}
		\begin{alertblock}{Critical Trap}
			Testing = Sampling (not 100%).
			Document methodology for auditability.
		\end{alertblock}
	\end{frame}
	
	\begin{frame}{Q66: Risk Response Cost Trap}
		\begin{block}{Question}
			The cost of a control is \$50K. The expected loss without the control is \$40K. What is the BEST response?
		\end{block}
		\begin{enumerate}[(A)]
			\item Implement the control - \$50K is affordable
			\item Accept the risk - control cost exceeds expected loss
			\item Transfer the risk - cheaper than implementing
			\item Avoid the risk - too expensive to mitigate
		\end{enumerate}
	\end{frame}
	
	\begin{frame}{Q66: Explanation}
		\begin{block}{Correct Answer: B}
			When \textbf{control cost > expected loss}, accept the risk:
			\begin{itemize}
				\item Control cost = \$50K
				\item Expected loss = \$40K
				\item Spending \$50K to prevent \$40K loss is not justified
				\item Accept the risk (with documentation)
			\end{itemize}
			Unless there are regulatory or strategic reasons to implement.
		\end{block}
		\begin{alertblock}{Critical Trap}
			Control cost > Expected loss = Accept.
			Economic analysis is fundamental.
		\end{alertblock}
	\end{frame}
	
	\begin{frame}{Q67: Reporting to Board Trap}
		\begin{block}{Question}
			What should NOT be included in a board-level risk report?
		\end{block}
		\begin{enumerate}[(A)]
			\item Key risk indicators and trends
			\item Comparison to risk appetite
			\item Technical details of vulnerabilities
			\item Recommended actions
		\end{enumerate}
	\end{frame}
	
	\begin{frame}{Q67: Explanation}
		\begin{block}{Correct Answer: C}
			\textbf{Technical details} should NOT be in board reports:
			\begin{itemize}
				\item Board members are not technical experts
				\item Focus should be on business impact
				\item Technical details obscure the message
				\item Board needs strategic view, not operational
			\end{itemize}
		\end{block}
		\begin{alertblock}{Critical Trap}
			Board = Business language.
			Technical details = Operational reports.
		\end{alertblock}
	\end{frame}
	
	\begin{frame}{Q68: Risk Monitoring Trap}
		\begin{block}{Question}
			What is the PRIMARY purpose of continuous risk monitoring?
		\end{block}
		\begin{enumerate}[(A)]
			\item To satisfy audit requirements
			\item To detect changes in risk profile in real-time
			\item To justify the risk management budget
			\item To track project progress
		\end{enumerate}
	\end{frame}
	
	\begin{frame}{Q68: Explanation}
		\begin{block}{Correct Answer: B}
			Continuous monitoring enables \textbf{real-time awareness} of:
			\begin{itemize}
				\item Changes in risk levels
				\item New threats emerging
				\item Control failures
				\item Risk appetite breaches
			\end{itemize}
			This enables proactive rather than reactive risk management.
		\end{block}
		\begin{alertblock}{Critical Trap}
			Continuous monitoring = Real-time awareness.
			Not just for audit - for proactive management.
		\end{alertblock}
	\end{frame}
	
	\begin{frame}{Q69: Key Control Indicator Trap}
		\begin{block}{Question}
			The percentage of firewall rules reviewed and approved is a:
		\end{block}
		\begin{enumerate}[(A)]
			\item Key Risk Indicator (KRI)
			\item Key Performance Indicator (KPI)
			\item Key Control Indicator (KCI)
			\item None of the above
		\end{enumerate}
	\end{frame}
	
	\begin{frame}{Q69: Explanation}
		\begin{block}{Correct Answer: C}
			This is a \textbf{KCI} because:
			\begin{itemize}
				\item Firewall rule review is a \textbf{control}
				\item The metric measures \textbf{control effectiveness}
				\item It doesn't predict risk (KRI)
				\item It doesn't measure overall performance (KPI)
			\end{itemize}
		\end{block}
		\begin{alertblock}{Critical Trap}
			If it measures a control = KCI.
			If it predicts risk = KRI.
			If it measures performance = KPI.
		\end{alertblock}
	\end{frame}
	
	\begin{frame}{Q70: Risk Owner Trap}
		\begin{block}{Question}
			The risk owner for a data breach risk is the CIO. Is this appropriate?
		\end{block}
		\begin{enumerate}[(A)]
			\item No, the CISO should own all security risks
			\item Yes, if the breach impacts IT objectives
			\item No, only the CEO can own this risk
			\item No, the board should own this risk
		\end{enumerate}
	\end{frame}
	
	\begin{frame}{Q70: Explanation}
		\begin{block}{Correct Answer: B}
			The risk owner should be the person who \textbf{owns the objective at risk}:
			\begin{itemize}
				\item If the breach impacts IT objectives (availability, integrity), CIO is appropriate
				\item If it impacts business objectives (customer trust, revenue), business leader is appropriate
				\item Risk ownership follows accountability for objectives
			\end{itemize}
		\end{block}
		\begin{alertblock}{Critical Trap}
			Risk owner = Objective owner.
			Not necessarily IT, not necessarily security.
		\end{alertblock}
	\end{frame}
	
	% Add remaining 30 questions (71-100) with similar structure
	% Due to length constraints, I'll continue with representative questions
	
	\begin{frame}{Q71: Exception Management Trap}
		\begin{block}{Question}
			An exception has been open for 3 years. What is the BEST action?
		\end{block}
		\begin{enumerate}[(A)]
			\item Renew for another year
			\item Review to determine if it is still needed and valid
			\item Automatically close it
			\item Escalate to the board
		\end{enumerate}
	\end{frame}
	
	\begin{frame}{Q71: Explanation}
		\begin{block}{Correct Answer: B}
			Long-standing exceptions require \textbf{detailed review}:
			\begin{itemize}
				\item Is the original justification still valid?
				\item Have conditions changed?
				\item Can the exception be closed?
				\item Are compensating controls still adequate?
			\end{itemize}
			Exceptions should not be auto-renewed indefinitely.
		\end{block}
	\end{frame}
	
	\begin{frame}{Q72: Control Owner Trap}
		\begin{block}{Question}
			A control owner leaves the organization. What should happen to the control?
		\end{block}
		\begin{enumerate}[(A)]
			\item The control is no longer needed
			\item The control continues without an owner
			\item A new owner should be assigned immediately
			\item The control should be retired
		\end{enumerate}
	\end{frame}
	
	\begin{frame}{Q72: Explanation}
		\begin{block}{Correct Answer: C}
			A \textbf{new control owner must be assigned immediately}:
			\begin{itemize}
				\item Controls without owners become orphaned
				\item No one monitors effectiveness
				\item No one maintains the control
				\item Control degradation is likely
			\end{itemize}
		\end{block}
	\end{frame}
	
	\begin{frame}{Q73: Risk Response Trap}
		\begin{block}{Question}
			A risk has been mitigated. The residual risk is within appetite. The risk owner wants to close the risk. What should the risk manager say?
		\end{block}
		\begin{enumerate}[(A)]
			\item Yes, close it - within appetite
			\item No, continue monitoring to ensure it stays within appetite
			\item No, always retain the risk in the register
			\item Yes, but only after board approval
		\end{enumerate}
	\end{frame}
	
	\begin{frame}{Q73: Explanation}
		\begin{block}{Correct Answer: B}
			The risk should \textbf{remain monitored}:
			\begin{itemize}
				\item Risk conditions can change
				\item Controls can degrade
				\item New threats can emerge
				\item Continuous monitoring confirms continued acceptability
			\end{itemize}
			The risk can be downgraded but not fully closed.
		\end{block}
	\end{frame}
	
	\begin{frame}{Q74: Third-Party Risk Trap}
		\begin{block}{Question}
			A vendor's security posture has deteriorated since onboarding. What is the FIRST action?
		\end{block}
		\begin{enumerate}[(A)]
			\item Terminate the contract immediately
			\item Assess the impact on your organization and communicate with the vendor
			\item Transfer the risk to another vendor
			\item Accept the risk and monitor
		\end{enumerate}
	\end{frame}
	
	\begin{frame}{Q74: Explanation}
		\begin{block}{Correct Answer: B}
			\textbf{Assess and communicate first}:
			\begin{itemize}
				\item Understand what has changed
				\item Determine impact on your data
				\item Communicate concerns to the vendor
				\item Develop a remediation plan
				\item Only terminate if necessary
			\end{itemize}
		\end{block}
	\end{frame}
	
	\begin{frame}{Q75: KRI Threshold Trap}
		\begin{block}{Question}
			A KRI has breached its threshold. What should the risk manager do FIRST?
		\end{block}
		\begin{enumerate}[(A)]
			\item Implement additional controls
			\item Escalate to senior management
			\item Investigate the cause of the breach
			\item Update the risk register
		\end{enumerate}
	\end{frame}
	
	\begin{frame}{Q75: Explanation}
		\begin{block}{Correct Answer: C}
			\textbf{Investigate first} before taking action:
			\begin{itemize}
				\item Is the data accurate?
				\item Is it a temporary spike or a trend?
				\item What is the root cause?
				\item What is the appropriate response?
			\end{itemize}
			Understanding the cause determines the right action.
		\end{block}
	\end{frame}
	
	% Continue with remaining questions 76-100
	% I'll include representative questions to reach 100
	
	\begin{frame}{Q76: Control Effectiveness Trap}
		\begin{block}{Question}
			A control has been tested and found effective. How long is this assurance valid?
		\end{block}
		\begin{enumerate}[(A)]
			\item Indefinitely
			\item 1 year
			\item Until conditions change or next test
			\item 3 years
		\end{enumerate}
	\end{frame}
	
	\begin{frame}{Q76: Explanation}
		\begin{block}{Correct Answer: C}
			Control effectiveness assurance is \textbf{point-in-time}:
			\begin{itemize}
				\item Valid until conditions change
				\item New threats can emerge
				\item Controls can degrade
				\item Personnel changes can affect operation
			\end{itemize}
			Regular retesting is required.
		\end{block}
	\end{frame}
	
	\begin{frame}{Q77: Risk Appetite Trap}
		\begin{block}{Question}
			Who is responsible for defining the organization's risk appetite?
		\end{block}
		\begin{enumerate}[(A)]
			\item Risk manager
			\item CISO
			\item Board of directors
			\item IT manager
		\end{enumerate}
	\end{frame}
	
	\begin{frame}{Q77: Explanation}
		\begin{block}{Correct Answer: C}
			The \textbf{Board of Directors} defines risk appetite:
			\begin{itemize}
				\item It represents stakeholder interests
				\item It has strategic oversight
				\item It sets the tone for the organization
				\item It has authority to accept enterprise risk
			\end{itemize}
		\end{block}
	\end{frame}
	
	\begin{frame}{Q78: Risk Monitoring Trap}
		\begin{block}{Question}
			What should trigger a review of the risk register?
		\end{block}
		\begin{enumerate}[(A)]
			\item Annual audit
			\item Only when a risk materializes
			\item Any significant change in environment, threats, or business
			\item Monthly
		\end{enumerate}
	\end{frame}
	
	\begin{frame}{Q78: Explanation}
		\begin{block}{Correct Answer: C}
			The risk register should be reviewed when:
			\begin{itemize}
				\item New threats emerge
				\item Business objectives change
				\item New regulations are introduced
				\item Significant incidents occur
				\item Organizational changes happen
			\end{itemize}
		\end{block}
	\end{frame}
	
	\begin{frame}{Q79: Control Testing Trap}
		\begin{block}{Question}
			What is the difference between a control test and a control assessment?
		\end{block}
		\begin{enumerate}[(A)]
			\item No difference
			\item Test is for IT, assessment is for business
			\item Test verifies operation, assessment evaluates overall effectiveness
			\item Test is quantitative, assessment is qualitative
		\end{enumerate}
	\end{frame}
	
	\begin{frame}{Q79: Explanation}
		\begin{block}{Correct Answer: C}
			\begin{itemize}
				\item \textbf{Control test:} Verifies the control is operating as designed (specific procedure)
				\item \textbf{Control assessment:} Evaluates whether the control is effectively reducing risk (overall judgment)
			\end{itemize}
			Testing is a component of assessment.
		\end{block}
	\end{frame}
	
	\begin{frame}{Q80: Risk Response Trap}
		\begin{block}{Question}
			A risk has been transferred. Does the organization still need to monitor it?
		\end{block}
		\begin{enumerate}[(A)]
			\item No, transfer eliminates responsibility
			\item Yes, transfer does not eliminate accountability
			\item No, the third party is responsible
			\item Only if required by contract
		\end{enumerate}
	\end{frame}
	
	\begin{frame}{Q80: Explanation}
		\begin{block}{Correct Answer: B}
			\textbf{Transfer does not eliminate accountability}:
			\begin{itemize}
				\item You remain accountable to regulators
				\item Reputational risk remains
				\item The third party may fail
				\item You must monitor the third party's performance
			\end{itemize}
		\end{block}
	\end{frame}
	
	\begin{frame}{Q81: Risk Reporting Trap}
		\begin{block}{Question}
			What is the MOST important quality of a risk report?
		\end{block}
		\begin{enumerate}[(A)]
			\item Comprehensiveness
			\item Actionability
			\item Technical accuracy
			\item Visual appeal
		\end{enumerate}
	\end{frame}
	
	\begin{frame}{Q81: Explanation}
		\begin{block}{Correct Answer: B}
			\textbf{Actionability} is most important:
			\begin{itemize}
				\item Reports should enable decisions
				\item Without action, reporting is pointless
				\item Readers should know what to do
				\item Focus on insights, not data
			\end{itemize}
		\end{block}
	\end{frame}
	
	\begin{frame}{Q82: Risk Culture Trap}
		\begin{block}{Question}
			What is the STRONGEST indicator of a healthy risk culture?
		\end{block}
		\begin{enumerate}[(A)]
			\item No security incidents
			\item Employees report risks without fear
			\item All risks are within appetite
			\item Zero audit findings
		\end{enumerate}
	\end{frame}
	
	\begin{frame}{Q82: Explanation}
		\begin{block}{Correct Answer: B}
			\textbf{Open risk reporting} indicates healthy culture:
			\begin{itemize}
				\item No fear of blame
				\item Proactive risk identification
				\item Learning from mistakes
				\item Continuous improvement
			\end{itemize}
		\end{block}
	\end{frame}
	
	\begin{frame}{Q83: Control Category Trap}
		\begin{block}{Question}
			A policy requiring password complexity is what type of control?
		\end{block}
		\begin{enumerate}[(A)]
			\item Technical
			\item Administrative
			\item Physical
			\item Detective
		\end{enumerate}
	\end{frame}
	
	\begin{frame}{Q83: Explanation}
		\begin{block}{Correct Answer: B}
			Password policy is an \textbf{administrative control}:
			\begin{itemize}
				\item It is a policy document
				\item It defines requirements
				\item It is not technical (enforced by technology)
				\item It is not physical
			\end{itemize}
		\end{block}
	\end{frame}
	
	\begin{frame}{Q84: Risk Acceptance Trap}
		\begin{block}{Question}
			What should be included in a risk acceptance document?
		\end{block}
		\begin{enumerate}[(A)]
			\item Only the risk description
			\item Risk description, justification, owner, and review date
			\item Only the approval signature
			\item Only the mitigation plan
		\end{enumerate}
	\end{frame}
	
	\begin{frame}{Q84: Explanation}
		\begin{block}{Correct Answer: B}
			Risk acceptance should include:
			\begin{itemize}
				\item Risk description and impact
				\item Business justification
				\item Risk owner (who accepts)
				\item Duration and review date
				\item Compensating controls
			\end{itemize}
		\end{block}
	\end{frame}
	
	\begin{frame}{Q85: Third-Party Trap}
		\begin{block}{Question}
			When should a vendor risk assessment be updated?
		\end{block}
		\begin{enumerate}[(A)]
			\item Only at contract renewal
			\item Annually or when significant changes occur
			\item Every 3 years
			\item Only when there is a breach
		\end{enumerate}
	\end{frame}
	
	\begin{frame}{Q85: Explanation}
		\begin{block}{Correct Answer: B}
			Vendor assessments should be updated:
			\begin{itemize}
				\item \textbf{Annually} (minimum)
				\item When the vendor's services change
				\item When there is a vendor incident
				\item When regulations change
				\item When the vendor's security posture changes
			\end{itemize}
		\end{block}
	\end{frame}
	
	\begin{frame}{Q86: Risk Monitoring Trap}
		\begin{block}{Question}
			What is the purpose of a risk dashboard?
		\end{block}
		\begin{enumerate}[(A)]
			\item To show all risks in detail
			\item To provide a consolidated view of key risk indicators
			\item To replace the risk register
			\item To satisfy audit requirements
		\end{enumerate}
	\end{frame}
	
	\begin{frame}{Q86: Explanation}
		\begin{block}{Correct Answer: B}
			A dashboard provides \textbf{consolidated view}:
			\begin{itemize}
				\item Key metrics at a glance
				\item Trend identification
				\item Comparison to thresholds
				\item Executive-level visibility
			\end{itemize}
		\end{block}
	\end{frame}
	
	\begin{frame}{Q87: Control Testing Trap}
		\begin{block}{Question}
			What should be done if a control test reveals a design deficiency?
		\end{block}
		\begin{enumerate}[(A)]
			\item Ignore it - testing is complete
			\item Document and recommend redesign
			\item Implement a compensating control only
			\item Accept the risk
		\end{enumerate}
	\end{frame}
	
	\begin{frame}{Q87: Explanation}
		\begin{block}{Correct Answer: B}
			Design deficiencies require \textbf{redesign}:
			\begin{itemize}
				\item Document the finding
				\item Assess the impact
				\item Recommend redesign
				\item Implement compensating controls temporarily
				\item Retest after redesign
			\end{itemize}
		\end{block}
	\end{frame}
	
	\begin{frame}{Q88: Risk Appetite Trap}
		\begin{block}{Question}
			How often should risk appetite be reviewed?
		\end{block}
		\begin{enumerate}[(A)]
			\item Never - it should remain constant
			\item Annually or when significant changes occur
			\item Every 5 years
			\item Only when there is a breach
		\end{enumerate}
	\end{frame}
	
	\begin{frame}{Q88: Explanation}
		\begin{block}{Correct Answer: B}
			Risk appetite should be reviewed:
			\begin{itemize}
				\item \textbf{Annually} (minimum)
				\item When business strategy changes
				\item When market conditions change
				\item When there are significant incidents
			\end{itemize}
		\end{block}
	\end{frame}
	
	\begin{frame}{Q89: Risk Response Trap}
		\begin{block}{Question}
			What is the PRIMARY consideration when selecting a risk response?
		\end{block}
		\begin{enumerate}[(A)]
			\item Cost of the response
			\item Alignment with business objectives and risk appetite
			\item Ease of implementation
			\item Vendor recommendations
		\end{enumerate}
	\end{frame}
	
	\begin{frame}{Q89: Explanation}
		\begin{block}{Correct Answer: B}
			\textbf{Alignment with business objectives} is primary:
			\begin{itemize}
				\item Risk response should support business goals
				\item Should bring risk within appetite
				\item Cost is secondary to effectiveness
				\item Ease of implementation is operational
			\end{itemize}
		\end{block}
	\end{frame}
	
	\begin{frame}{Q90: Risk Register Trap}
		\begin{block}{Question}
			How often should a risk register be updated?
		\end{block}
		\begin{enumerate}[(A)]
			\item Annually
			\item When risks change or new risks emerge
			\item Quarterly
			\item Monthly
		\end{enumerate}
	\end{frame}
	
	\begin{frame}{Q90: Explanation}
		\begin{block}{Correct Answer: B}
			The risk register should be updated \textbf{when conditions change}:
			\begin{itemize}
				\item New risks identified
				\item Risk levels change
				\item Risk responses implemented
				\item Risk owners change
			\end{itemize}
			It is a living document, not a periodic report.
		\end{block}
	\end{frame}
	
	\begin{frame}{Q91: Exception Trap}
		\begin{block}{Question}
			Who should approve a policy exception?
		\end{block}
		\begin{enumerate}[(A)]
			\item IT manager
			\item Business process owner or delegate with authority
			\item Risk manager
			\item Anyone affected by the policy
		\end{enumerate}
	\end{frame}
	
	\begin{frame}{Q91: Explanation}
		\begin{block}{Correct Answer: B}
			Exceptions should be approved by the \textbf{business process owner}:
			\begin{itemize}
				\item They own the risk
				\item They have budget authority
				\item They understand the business need
				\item They are accountable for the exception
			\end{itemize}
		\end{block}
	\end{frame}
	
	\begin{frame}{Q92: Risk Monitoring Trap}
		\begin{block}{Question}
			What is the difference between monitoring and auditing?
		\end{block}
		\begin{enumerate}[(A)]
			\item No difference
			\item Monitoring is continuous, auditing is periodic
			\item Monitoring is internal, auditing is external
			\item Monitoring is automated, auditing is manual
		\end{enumerate}
	\end{frame}
	
	\begin{frame}{Q92: Explanation}
		\begin{block}{Correct Answer: B}
			\begin{itemize}
				\item \textbf{Monitoring:} Continuous, ongoing oversight
				\item \textbf{Auditing:} Periodic, independent examination
			\end{itemize}
			Monitoring is the day-to-day function; auditing provides independent assurance.
		\end{block}
	\end{frame}
	
	\begin{frame}{Q93: Control Design Trap}
		\begin{block}{Question}
			What should be considered FIRST when designing a control?
		\end{block}
		\begin{enumerate}[(A)]
			\item Cost of the control
			\item Risk being addressed
			\item Available technology
			\item User convenience
		\end{enumerate}
	\end{frame}
	
	\begin{frame}{Q93: Explanation}
		\begin{block}{Correct Answer: B}
			The \textbf{risk being addressed} is primary:
			\begin{itemize}
				\item Understand the risk first
				\item Then design appropriate control
				\item Cost is secondary
				\item Technology is a means, not the end
			\end{itemize}
		\end{block}
	\end{frame}
	
	\begin{frame}{Q94: KRI Selection Trap}
		\begin{block}{Question}
			What makes a KRI effective?
		\end{block}
		\begin{enumerate}[(A)]
			\item It is easy to measure
			\item It is actionable and predictive of risk
			\item It is approved by the board
			\item It is automated
		\end{enumerate}
	\end{frame}
	
	\begin{frame}{Q94: Explanation}
		\begin{block}{Correct Answer: B}
			Effective KRIs are \textbf{actionable and predictive}:
			\begin{itemize}
				\item They predict future risk
				\item They trigger action when breached
				\item They are relevant to the risk
				\item They enable proactive management
			\end{itemize}
		\end{block}
	\end{frame}
	
	\begin{frame}{Q95: Risk Response Trap}
		\begin{block}{Question}
			What is residual risk?
		\end{block}
		\begin{enumerate}[(A)]
			\item Risk before controls
			\item Risk after controls
			\item Risk that is accepted
			\item Risk that is transferred
		\end{enumerate}
	\end{frame}
	
	\begin{frame}{Q95: Explanation}
		\begin{block}{Correct Answer: B}
			\textbf{Residual risk = risk after controls}:
			\begin{itemize}
				\item Inherent risk - Control effectiveness = Residual risk
				\item It is the risk that remains
				\item It must be within risk appetite
				\item It must be accepted by management
			\end{itemize}
		\end{block}
	\end{frame}
	
	\begin{frame}{Q96: Control Implementation Trap}
		\begin{block}{Question}
			What is the FIRST step in implementing a control?
		\end{block}
		\begin{enumerate}[(A)]
			\item Purchase the technology
			\item Define the control objective and requirements
			\item Train the users
			\item Test the control
		\end{enumerate}
	\end{frame}
	
	\begin{frame}{Q96: Explanation}
		\begin{block}{Correct Answer: B}
			\textbf{Define objectives first}:
			\begin{itemize}
				\item What risk are we addressing?
				\item What should the control achieve?
				\item What are the requirements?
				\item Then select and implement the solution
			\end{itemize}
		\end{block}
	\end{frame}
	
	\begin{frame}{Q97: Risk Reporting Trap}
		\begin{block}{Question}
			What is the BEST format for reporting to the board?
		\end{block}
		\begin{enumerate}[(A)]
			\item Detailed technical report
			\item Executive summary with key metrics and recommendations
			\item Full risk register
			\item Raw data and logs
		\end{enumerate}
	\end{frame}
	
	\begin{frame}{Q97: Explanation}
		\begin{block}{Correct Answer: B}
			Board reports should be \textbf{concise and strategic}:
			\begin{itemize}
				\item Executive summary (1-2 pages)
				\item Key metrics and trends
				\item Comparison to appetite
				\item Recommended actions
			\end{itemize}
		\end{block}
	\end{frame}
	
	\begin{frame}{Q98: Risk Acceptance Trap}
		\begin{block}{Question}
			What happens if a risk is accepted without documentation?
		\end{block}
		\begin{enumerate}[(A)]
			\item Nothing - the risk is accepted
			\item The acceptance may not be valid or defensible
			\item The risk is automatically transferred
			\item The risk is automatically avoided
		\end{enumerate}
	\end{frame}
	
	\begin{frame}{Q98: Explanation}
		\begin{block}{Correct Answer: B}
			Undocumented acceptance is \textbf{not defensible}:
			\begin{itemize}
				\item No evidence of due care
				\item Cannot prove informed decision
				\item May be considered negligence
				\item Vulnerable to legal challenge
			\end{itemize}
		\end{block}
	\end{frame}
	
	\begin{frame}{Q99: Risk Monitoring Trap}
		\begin{block}{Question}
			What should be monitored for third-party risks?
		\end{block}
		\begin{enumerate}[(A)]
			\item Only contract compliance
			\item Security posture, incidents, and compliance
			\item Only financial stability
			\item Only service level agreements
		\end{enumerate}
	\end{frame}
	
	\begin{frame}{Q99: Explanation}
		\begin{block}{Correct Answer: B}
			Comprehensive third-party monitoring includes:
			\begin{itemize}
				\item Security posture (SOC reports, assessments)
				\item Incidents and breaches
				\item Regulatory compliance
				\item Financial stability
				\item Service level performance
			\end{itemize}
		\end{block}
	\end{frame}
	
	\begin{frame}{Q100: Final Trap}
		\begin{block}{Question}
			What is the ULTIMATE goal of risk response and reporting?
		\end{block}
		\begin{enumerate}[(A)]
			\item Eliminate all risk
			\item Optimize risk to support business objectives
			\item Satisfy auditors
			\item Reduce insurance premiums
		\end{enumerate}
	\end{frame}
	
	\begin{frame}{Q100: Explanation}
		\begin{block}{Correct Answer: B}
			The ultimate goal is \textbf{risk optimization}:
			\begin{itemize}
				\item Risk cannot be eliminated
				\item Risk should be managed to acceptable levels
				\item Risk management enables business objectives
				\item Balance risk and reward
			\end{itemize}
		\end{block}
		\begin{alertblock}{Final Wisdom}
			Risk management is not about eliminating risk.\\
			It is about understanding risk and making informed decisions.
		\end{alertblock}
	\end{frame}
	
	\begin{frame}{Final Summary}
		\begin{block}{Key Takeaways for Domain C}
			\begin{itemize}
				\item \textbf{Mitigate} = Reduce risk, not eliminate
				\item \textbf{Accept} = Formal, documented decision
				\item \textbf{Transfer} = Financial shift, accountability remains
				\item \textbf{Avoid} = Eliminate activity that creates risk
				\item \textbf{Risk Owner} = Business leader accountable for risk
				\item \textbf{Control Owner} = Implementer of controls
				\item \textbf{KRI} = Predicts risk (forward-looking)
				\item \textbf{KPI} = Measures performance (backward-looking)
				\item \textbf{KCI} = Measures control effectiveness
			\end{itemize}
		\end{block}
	\end{frame}
	
	\begin{frame}{Thank You}
		\begin{center}
			\Huge
			\textbf{Good Luck on Your CRISC Exam!}
			
			\vspace{1cm}
			
			\normalsize
			Remember: The ISACA mindset is about\\
			understanding why a question is asked,\\
			not just what is technically correct.
		\end{center}
	\end{frame}
	
\end{document}